% from ega3-1-fr.tex
\oldpage[0\textsubscript{III}]{15}
La condition est évidemment suffisante, tout ensemble ouvert non vide étant
dense dans $X$. Inversement, soit
$E=\bigcup_{i=1}^n(U_i\cap F_i)$ une partie constructible de $X$, les $U_i$
étant ouverts non vides et les $F_i$ fermés dans $X$~; on a donc
$\overline{E}\subset\bigcup_i F_i$. Par suite, si $\overline{E}=X$, $X$ est
égal à l'un des $F_i$, donc $E\supset U_i$, ce qui achève la démonstration.

Lorsque $X$ admet un point générique $x$
(\hyperref[0.2.1.2-fr]{$0_{\mathrm I}$, 2.1.2}), la condition de
(\hyperref[0.9.2.2-fr]{9.2.2}) équivaut à la relation $x\in E$.

\begin{proposition}[9.2.3]
\label{0.9.2.3-fr}
Soit $X$ un espace noethérien. Pour qu'une partie $E$ de $X$ soit
constructible, il faut et il suffit que, pour toute partie fermée
irréductible $Y$ de $X$, $E\cap Y$ soit rare dans $Y$ ou contienne une partie
ouverte non vide de $Y$.
\end{proposition}

La nécessité de la condition provient de ce que $E\cap Y$ doit être une
partie constructible de $Y$ et de (\hyperref[0.9.2.2-fr]{9.2.2}), car une
partie non dense de $Y$ est nécessairement rare dans l'espace irréductible
$Y$ (\hyperref[0.2.1.1-fr]{$0_{\mathrm I}$, 2.1.1}). Pour prouver que la
condition est suffisante, appliquons le principe de récurrence noethérienne
(\hyperref[0.2.2.2-fr]{$0_{\mathrm I}$, 2.2.2}) à l'ensemble
$\mathfrak{F}$ des parties fermées $Y$ de $X$ telles que $Y\cap E$ soit
constructible (par rapport à $Y$ ou par rapport à $X$, ce qui revient au
même)~: on peut donc supposer que pour toute partie fermée $Y\ne X$ de $X$,
$E\cap Y$ est constructible. Supposons d'abord que $X$ ne soit pas
irréductible, et soient $X_i$ ($1\leq i\leq m$) ses composantes
irréductibles, nécessairement en nombre fini
(\hyperref[0.2.2.5-fr]{$0_{\mathrm I}$, 2.2.5})~; par hypothèse, les
$E\cap X_i$ sont constructibles, donc aussi leur réunion $E$. Supposons
ensuite que $X$ soit irréductible~; alors, par hypothèse, ou bien $E$ est
rare, donc $\overline{E}\ne X$ et $E=E\cap\overline{E}$ est constructible~;
ou bien $E$ contient un ouvert non vide $U$ de $E$, donc est réunion de $U$
et de $E\cap(X\setminus U)$~; mais $X\setminus U$ est un ensemble fermé
distinct de $X$, donc $E\cap(X\setminus U)$ est constructible~; $E$
lui-même est par suite constructible, ce qui achève la démonstration.

\begin{corollary}[9.2.4]
\label{0.9.2.4-fr}
Soient $X$ un espace noethérien, $(E_\alpha)$ une famille filtrante
croissante de parties constructibles de $X$, telle que~:
\begin{enumerate}
  \item[$1^\circ$] $X$ est réunion de la famille $(E_\alpha)$.
  \item[$2^\circ$] Toute partie fermée irréductible de $X$ est contenue dans
    l'adhérence d'un des $E_\alpha$.
\end{enumerate}
Alors il existe un indice $\alpha$ tel que $X=E_\alpha$.

Lorsque toute partie fermée irréductible de $X$ admet un point générique,
l'hypothèse $2^\circ$ peut être supprimée.
\end{corollary}

Appliquons le principe de récurrence noethérienne
(\hyperref[0.2.2.2-fr]{$0_{\mathrm I}$, 2.2.2}) à l'ensemble
$\mathfrak{M}$ des parties fermées de $X$ contenues dans l'un des $E_\alpha$
au moins~; on peut donc supposer que toute partie fermée $Y\ne X$ de $X$ est
contenue dans un des $E_\alpha$. La proposition est évidente si $X$ n'est pas
irréductible, car chacune des composantes irréductibles $X_i$ de $X$
($1\leq i\leq m$) est contenue dans un $E_{\alpha_i}$, et il existe un
$E_\alpha$ contenant tous les $E_{\alpha_i}$. Supposons donc $X$
irréductible. Par hypothèse, il existe $\beta$ tel que
$X=\overline{E_\beta}$, donc (\hyperref[0.9.2.2-fr]{9.2.2}) $E_\beta$
contient un ouvert non vide $U$ de $X$. Mais alors l'ensemble fermé
$X\setminus U$ est contenu dans un $E_\gamma$, et il suffit de prendre
$E_\alpha$ contenant $E_\beta$ et $E_\gamma$. Lorsque toute partie fermée
irréductible $Y$ de $X$
